\addcontentsline{toc}{section}{References}
\begin{thebibliography}{99}
\small

\bibitem{ciocan} B. I. Ciocan et al. (2026),
``MUSE-DARK III: The evolution of the radial acceleration relation at
intermediate redshifts,'' \emph{Astronomy \& Astrophysics} 709, L16.
\href{https://arxiv.org/abs/2604.22613}{arXiv:2604.22613}.
The coefficient comparison uses Appendix E's self-consistent inference.

\bibitem{unruh} W. G. Unruh (1976), ``Notes on black-hole evaporation,''
\emph{Physical Review D} 14, 870.
\href{https://doi.org/10.1103/PhysRevD.14.870}{doi:10.1103/PhysRevD.14.870}.

\bibitem{gibbons} G. W. Gibbons and S. W. Hawking (1977),
``Cosmological event horizons, thermodynamics, and particle creation,''
\emph{Physical Review D} 15, 2738.
\href{https://doi.org/10.1103/PhysRevD.15.2738}{doi:10.1103/PhysRevD.15.2738}.

\bibitem{warren} H. T. Warren et al. (2025), ``Estimate of Current Mass
of the Large Magellanic Cloud from the Orphan-Chenab Tidal Stream.''
\href{https://arxiv.org/abs/2508.01086}{arXiv:2508.01086}.

\bibitem{brooks} R. A. N. Brooks et al. (2026), ``Quantifying the
Milky Way, LMC and their interaction using all-sky kinematics of outer
halo stars.'' \href{https://arxiv.org/abs/2603.08788}{arXiv:2603.08788}.

\bibitem{lelli} F. Lelli et al. (2015), ``Gas dynamics in tidal dwarf
galaxies: disc formation at $z=0$.''
\href{https://arxiv.org/abs/1509.05404}{arXiv:1509.05404}.

\bibitem{wolf} J. Wolf et al. (2010), ``Accurate masses for
dispersion-supported galaxies,'' \emph{MNRAS} 406, 1220.
\href{https://arxiv.org/abs/0908.2995}{arXiv:0908.2995}.

\bibitem{moster} B. P. Moster, T. Naab, and S. D. M. White (2013),
``Galactic star formation and accretion histories from matching galaxies
to dark matter haloes.''
\href{https://arxiv.org/abs/1205.5807}{arXiv:1205.5807}.

\bibitem{dutton} A. A. Dutton and A. V. Macci\`o (2014),
``Cold dark matter haloes in the Planck era: evolution of structural
parameters for Einasto and NFW profiles.''
\href{https://arxiv.org/abs/1402.7073}{arXiv:1402.7073}.

\bibitem{atzberger} K. R. Atzberger et al. (2026),
``The DECam MAGIC Survey: Uncovering the Tidal Tails of the Crater II
Dwarf Galaxy.''
\href{https://arxiv.org/abs/2602.21283}{arXiv:2602.21283}.

\bibitem{limberg} G. Limberg et al. (2025),
``$S^5$: Tidal Disruption in Crater 2 and Formation of Diffuse Dwarf
Galaxies in the Local Group.''
\href{https://arxiv.org/abs/2512.02177}{arXiv:2512.02177}.

\bibitem{ji} A. P. Ji et al. (2021),
``Kinematics of Antlia 2 and Crater 2 from The Southern Stellar Stream
Spectroscopic Survey ($S^5$).''
\href{https://arxiv.org/abs/2106.12656}{arXiv:2106.12656}.

\bibitem{desi} DESI Collaboration (2025),
``DESI DR2 Results II: Measurements of Baryon Acoustic Oscillations and
Cosmological Constraints.''
\href{https://arxiv.org/abs/2503.14738}{arXiv:2503.14738}.

\bibitem{planck} Planck Collaboration (2020),
``Planck 2018 results. VI. Cosmological parameters.''
\href{https://arxiv.org/abs/1807.06209}{arXiv:1807.06209}.

\bibitem{desrecal} B. Popovic et al. (2025),
``The Dark Energy Survey Supernova Program: A Reanalysis Of Cosmology
Results And Evidence For Evolving Dark Energy With An Updated Type Ia
Supernova Calibration.''
\href{https://arxiv.org/abs/2511.07517}{arXiv:2511.07517}.

\bibitem{desrepo} DES Science, public supernova data and covariance
repository. \href{https://github.com/des-science/DES-SN5YR}{DES-SN5YR on
GitHub}. The archived data version used by the calculations should be
preserved when reproducing them.

\bibitem{releasecode} J. Chinitz, public dark-energy reproduction release.
\href{https://github.com/jakechinitz/dark-energy/releases/tag/Dark-Energy}
{Dark-Energy release on GitHub}. The current mixing extension is
documented by local sources C1--C3.


\bibitem{xcopnew} S. Ettori et al., ``Hydrostatic mass profiles in X-COP galaxy clusters,'' A\&A 621, A39 (2019).
\href{https://arxiv.org/abs/1805.00035}{arXiv:1805.00035};
\href{https://dominiqueeckert.wixsite.com/xcop/data}{public X-COP data}.
\bibitem{herbonnetnew} R. Herbonnet et al., ``CCCP and MENeaCS: (updated) weak-lensing masses for 100 galaxy clusters,'' MNRAS 497, 4684 (2020), Table A2.
\href{https://arxiv.org/abs/1912.04414}{arXiv:1912.04414}.
\bibitem{bulletnew} G. Rihtar\v{s}i\v{c} et al., ``Mapping dark matter in the Bullet Cluster using JWST imaging and spectroscopy'' (2026).
\href{https://arxiv.org/html/2601.22245v2}{arXiv:2601.22245v2};
\href{https://niriss.github.io/lensing.html}{CANUCS Bullet v1 maps}.
\bibitem{sarkarnew} A. Sarkar et al., ``XRISM Observations of Abell 1795: Evidence for Low Turbulence and Resonant Scattering'' (2026).
\href{https://arxiv.org/abs/2606.08097}{arXiv:2606.08097}.
\bibitem{xrismnew} XRISM Collaboration, ``Constraining gas motion and non-thermal pressure beyond the core of the Abell 2029 galaxy cluster with XRISM'' (2025).
\href{https://arxiv.org/abs/2505.06533}{arXiv:2505.06533}.

\end{thebibliography}
